\documentclass[preview]{standalone}

\usepackage{amsmath}
\usepackage{amssymb}
\usepackage{stellar}
\usepackage{definitions}

\begin{document}

\id{module-generators-exercises}
\genpage

\section{Exercises}

\begin{snippetexercise}{modules-free-and-nonfree-exercise}{Free and nonfree modules}
    Let \(K\) be a \field.
    \begin{enumerate}
        \item Let
        \[
            A=
            \left\{
                \begin{pmatrix}a&b\\b&a\end{pmatrix}
                \suchthat a,b\in K
            \right\}
            \subringleq\matrices_2(K),
            \qquad
            M=K^2,
        \]
        where \(M\) is an \(A\)-\module under matrix multiplication.
        Show that \(M\) is a \freemodule over \(A\).

        \item Let
        \[
            A=K[X],
            \qquad
            I=(X^2-1),
            \qquad
            M=A/I.
        \]
        Show that \(M\) is finitely generated as an \(A\)-module but is
        not free.
    \end{enumerate}
\end{snippetexercise}

\begin{snippetsolution}{modules-free-and-nonfree-solution}{Free and nonfree modules}
    For the first module, put
    \[
        e_1=\begin{pmatrix}1\\0\end{pmatrix}.
    \]
    Every vector satisfies
    \[
        \begin{pmatrix}x\\y\end{pmatrix}
        =
        \begin{pmatrix}x&y\\y&x\end{pmatrix}e_1,
    \]
    so \(e_1\) generates \(M\). Moreover,
    \[
        \begin{pmatrix}a&b\\b&a\end{pmatrix}e_1=0_M
        \quad\Longrightarrow\quad
        \begin{pmatrix}a\\b\end{pmatrix}=0_M
        \quad\Longrightarrow\quad
        a=b=0.
    \]
    Thus \(\{e_1\}\) is an \(A\)-basis and \(M\) is free of rank \(1\).

    For the second module, every element \(f+I\) equals
    \(f(1+I)\), so \(1+I\) generates \(M\). However,
    \[
        (X^2-1)(1+I)=0_M,
        \qquad
        X^2-1\neq0,
        \qquad
        1+I\neq0_M.
    \]
    Hence \(M\) has nonzero torsion. A free module over the
    \integraldomain \(K[X]\) is torsion-free, so \(M\) is not free.
\end{snippetsolution}

\end{document}
